\begin{center}
    \section*{\fontsize{20}{20}\selectfont Chapter 6}
\end{center}
\vspace{10mm}

\section{Conclusion}

This thesis demonstrates that graph neural networks constitute a viable and computationally efficient approach to brain tumour segmentation, challenging the prevailing assumption that volumetric 3D convolutions are necessary for competitive medical image analysis. The proposed framework achieves \textbf{91.41\% Dice coefficient} through soft-voting ensemble aggregation of five cross-validation models, evaluated on a sealed 251-patient held-out set. Individual fold performance averages \textbf{90.02\% $\pm$ 0.74\% Dice} across five folds with 720/80/200 train/validation/test patient splits — demonstrating consistent generalisation without overfitting. The framework delivers \textbf{6.9$\times$ faster inference} (74\,ms per patient on GPU with pre-built graphs; 1.47\,s end-to-end including fast superpixel construction) and \textbf{155$\times$ fewer parameters} (439K vs.\ 68M) compared to a standard 3D U-Net baseline, with a peak GPU memory footprint of only \textbf{11\,MB}. These efficiency gains enable deployment on consumer-grade hardware (NVIDIA RTX 2060, 6\,GB VRAM), representing a concrete step towards democratising AI-assisted diagnostics in resource-constrained clinical settings.

Zero-shot cross-dataset evaluation on BraTS 2023 (1,245 patients) yields \textbf{89.21\% Dice} — a 2.20\% generalisation gap consistent with expected domain shift between dataset editions — confirming that the model is not overfit to a single dataset distribution and is capable of cross-institutional generalisation.

The principal contributions of this work are fivefold. First, a superpixel-based graph construction pipeline that exploits inherent tumour sparsity to achieve at least 890$\times$ spatial dimensionality reduction (measured per-patient compression $\approx$2,284$\times$; the conservative 890$\times$ figure uses the theoretical SLIC target of 200 superpixels per slice at an assumed 30\% brain fraction — see Chapter 3 Eq.~\ref{eq:compression_890}--\ref{eq:compression_actual} for full derivation) while preserving tumour-discriminative features across four MRI modalities (T1, T1CE, T2, FLAIR). Second, rigorous evaluation via patient-level stratified 5-fold cross-validation with a sealed held-out test set and 15 automated integrity checks, establishing zero data leakage and full reproducibility. Third, empirical validation through systematic ablation studies (depth, width, architecture) that the 5-layer, 256-dim GraphSAGE is the Pareto-optimal design under efficiency constraints. Fourth, zero-shot cross-dataset generalisation to BraTS 2023, quantifying a 2.20\% domain shift gap. Fifth, qualitative failure case analysis identifying complete-miss failures on atypical tumours as the primary error mode and providing concrete directions for future robustness improvements.

Returning to the central research question — \textit{Can a graph-based neural network achieve competitive segmentation accuracy ($>$90\% Dice) while providing efficiency gains sufficient for resource-constrained clinical deployment?} — this thesis provides a substantively affirmative answer. The 91.41\% ensemble Dice, combined with 74\,ms inference time, 11\,MB memory footprint, and 439K parameter count, demonstrates that competitive accuracy and practical efficiency are not mutually exclusive. As healthcare AI scales to serve diverse global populations, efficiency-aware architectures of this kind will be critical for ensuring that advanced diagnostic tools remain accessible regardless of hardware infrastructure.


\subsection{Limitations}

Several limitations of this work are acknowledged. First, graph construction (SLIC superpixel segmentation and feature extraction) adds approximately 1.5 seconds per patient as a preprocessing step; while negligible when graphs are cached, this adds latency in raw-scan real-time scenarios. Second, the 15 handcrafted node features require domain expertise in MRI physics and tumour biology; they may not generalise without modification to tumour types outside the glioma spectrum represented in BraTS. Third, this work addresses binary tumour segmentation only, foregoing the clinically relevant sub-region delineation (Enhancing Tumour, Tumour Core, Whole Tumour) evaluated in standard BraTS multi-class benchmarks. Fourth, while BraTS 2023 zero-shot evaluation provides preliminary evidence of cross-dataset generalisation, performance on entirely different scanner protocols, patient populations, or tumour subtypes (e.g.\ meningioma, metastases) remains untested. Fifth, a 1.6--2.4\% Dice gap remains below the best transformer-based methods on the Whole Tumour metric. Sixth, clinical integration barriers — including regulatory approval (FDA 510(k) / CE marking), hospital IT infrastructure requirements, and prospective radiologist validation studies — remain unaddressed and constitute necessary steps before clinical translation.


\subsection{Future Work and Directions}

Several directions emerge naturally from this work. Architecturally, incorporating attention mechanisms (e.g.\ graph attention layers) or hierarchical multi-scale graph representations could close the remaining gap with transformer-based methods without sacrificing efficiency. Ablation results (single-fold, Fold 0 protocol) suggest that widening hidden dimensions from 256 to 512 yields a +4.75\% improvement (88.78\% vs.\ 84.03\%) at the cost of 3.9$\times$ more parameters — a viable trade-off in settings with more computational budget. A multi-task extension enabling simultaneous binary and multi-class segmentation (ET, TC, WT) would bring the approach to full BraTS benchmark comparability and significantly increase clinical utility.

On the preprocessing side, GPU-accelerated superpixel segmentation or differentiable clustering approaches could reduce the graph construction overhead, enabling true end-to-end real-time processing. Addressing the dominant complete-miss failure mode — through adaptive superpixel counts for small lesions, data augmentation for atypical tumour appearances, or test-time threshold calibration — would improve robustness across the full patient distribution. Federated learning extensions would leverage the lightweight model architecture to enable privacy-preserving multi-institutional training without sharing patient data.

For clinical translation, the immediate priorities are prospective radiologist evaluation studies, multi-institutional validation on diverse scanner hardware, and regulatory submission preparation. Application of the same graph-based pipeline to other anatomical regions (prostate, liver, lung) and to longitudinal tumour tracking for treatment response assessment are further avenues of investigation.

This thesis establishes graph neural networks as a compelling paradigm for efficient brain tumour segmentation, demonstrating that competitive medical image analysis does not require massive volumetric models. The combination of strong accuracy (91.41\% ensemble Dice), exceptional efficiency (74\,ms, 11\,MB, 439K parameters), and demonstrated cross-dataset generalisation (89.21\% on BraTS 2023 zero-shot) positions graph-based segmentation as a practical tool for the democratisation of AI-assisted diagnostics across all levels of healthcare infrastructure.
